% from ega2-1-fr.tex
\oldpage[II]{176}
  un $Y$-morphisme $v:C\to C''$ (avec
  $C''=\operatorname{Spec}(\mathcal{S}'')$) tel que
  $\mathcal{A}(v)=\varphi$, et comme $\varphi$ est un homomorphisme
  d'Algèbres graduées, on déduit de $\varphi$ un $v$-morphisme d'Algèbres
  graduées
  $\psi:\mathcal{S}''^{\natural}\to\mathcal{S}^{\natural}$
  (\hyperref[II.1.4.1-fr]{1.4.1}). En outre, il résulte aisément de
  (\hyperref[II.8.2.10-fr]{8.2.10, (iv)}) et de l'hypothèse sur $\varphi$,
  que $\operatorname{Proj}(\psi)$ est partout défini. Enfin, compte tenu de
  (\hyperref[II.3.5.6.1-fr]{3.5.6.1}), on a un $u$-morphisme canonique
  $\mathcal{S}_{X''}\to\mathcal{S}_X$, d'où
  (\hyperref[II.1.5.6-fr]{1.5.6}) un morphisme $w:C_{X''}\to C_X$. Cela
  étant, le diagramme
  \begin{equation}
  \label{II.8.7.4.2-fr}
    \xymatrix{
      C_{X''}\ar[r]^{\sim}\ar[d]_{w} &
      \operatorname{Proj}(\mathcal{S}^{\natural})
        \ar[d]^{\operatorname{Proj}(\psi)}\\
      C_X\ar[r]^{\sim} & \operatorname{Proj}(\mathcal{S}^{\natural})
    }
  \tag{8.7.4.2}
  \end{equation}
  est commutatif, comme on le vérifie aussitôt en se ramenant au cas où $Y$
  est affine.

  \item[\rm{(ii)}] Notons qu'en vertu de
  (\hyperref[II.8.7.2.5-fr]{8.7.2.5}) et
  (\hyperref[II.8.7.2.6-fr]{8.7.2.6}), on a
  $\mathcal{I}_1^m\subset\mathcal{I}_m\subset\mathcal{I}_1$ pour tout $m>0$.
  Or, par définition, $\mathcal{I}_1=(\mathcal{S}_+)^{\sim}$, donc
  $\mathcal{I}_1$ définit dans $C$ le sous-préschéma fermé
  $\varepsilon(Y)$ ((\hyperref[II.1.4.10-fr]{1.4.10}) et
  (\hyperref[II.8.3.2-fr]{8.3.2}))~; on en conclut que pour tout $m>0$,
  \emph{le support de $\mathcal{O}_C/\mathcal{I}_m$ est contenu dans l'espace
  sous-jacent au préschéma des sommets $\varepsilon(Y)$}~; dans l'image
  réciproque du cône épointé affine $E$, le morphisme structural
  $\operatorname{Proj}(\mathcal{S}^{\natural})\to C$ se réduit donc à un
  \emph{isomorphisme} (comme il résulte d'ailleurs de
  (\hyperref[II.8.7.3-fr]{8.7.3}) et
  (\hyperref[II.8.7.1-fr]{8.7.1})). En outre, identifiant canoniquement $C$ à
  un ouvert de $\widehat{C}$ (\hyperref[II.8.3.3-fr]{8.3.3}), on peut
  évidemment prolonger les Idéaux $\mathcal{I}_m$ de $\mathcal{O}_C$ en des
  Idéaux $\mathcal{J}_m$ de $\mathcal{O}_{\widehat{C}}$, par la condition de
  coïncider avec $\mathcal{O}_{\widehat{C}}$ dans l'ouvert $\widehat{E}$ de
  $\widehat{C}$. Si on pose
  $\mathcal{T}=\bigoplus_{n\geqslant 0}\mathcal{J}_m$, qui est une
  $\mathcal{O}_{\widehat{C}}$-Algèbre graduée quasi-cohérente, on peut
  prolonger l'isomorphisme (\hyperref[II.8.7.3.1-fr]{8.7.3.1}) en un
  $\widehat{C}$-isomorphisme
  \begin{equation}
  \label{II.8.7.4.3-fr}
    \widehat{C}_X\xrightarrow{\sim}\operatorname{Proj}(\mathcal{T}).
  \tag{8.7.4.3}
  \end{equation}

  En effet, au-dessus de $\widehat{E}$, il résulte de ce qui précède que
  $\operatorname{Proj}(\mathcal{T})$ s'identifie canoniquement à
  $\widehat{E}$, et on définit donc l'isomorphisme
  (\hyperref[II.8.7.4.3-fr]{8.7.4.3}) au-dessus de $\widehat{E}$ en lui
  imposant de coïncider avec l'isomorphisme canonique
  $\widehat{E}_X\to\widehat{E}$ (\hyperref[II.8.6.2-fr]{8.6.2})~; il est clair
  qu'alors cet isomorphisme et (\hyperref[II.8.7.3.1-fr]{8.7.3.1}) coïncident
  au-dessus de $E$.
\end{enumerate}
\end{remark}

\begin{corollary}[8.7.5]
\label{II.8.7.5-fr}
Supposons qu'il existe $n_0>0$ tel que
\begin{equation}
\label{II.8.7.5.1-fr}
  \mathcal{S}_{n+1}=\mathcal{S}_1\mathcal{S}_n
  \qquad\text{pour }n\geqslant n_0.
\tag{8.7.5.1}
\end{equation}
Alors le sous-préschéma des sommets de $C_X$ (isomorphe à $X$) est l'image
réciproque par le morphisme canonique $r:C_X\to C$ du sous-préschéma des
sommets de $C$ (isomorphe à $Y$). Inversement, si cette propriété a lieu et
si on suppose de plus $Y$ noethérien et $\mathcal{S}$ de type fini, il existe
$n_0>0$ tel que l'on ait (\hyperref[II.8.7.5.1-fr]{8.7.5.1}).
\end{corollary}

\begin{proof}
La première assertion étant locale sur $Y$, on peut pour l'établir supposer
$Y=\operatorname{Spec}(A)$ affine, et $\mathcal{S}=\widetilde{S}$, où $S$ est
une $A$-algèbre graduée à degrés positifs. Elle résulte alors de
(\hyperref[II.8.2.12-fr]{8.2.12}), car on a
$\operatorname{Proj}(S^{\natural}\otimes_S S_0)=
C_X\times_C\varepsilon(Y)$ (en vertu de l'identification
(\hyperref[II.8.7.3.1-fr]{8.7.3.1})), c'est-à-dire que ce préschéma est
l'image réciproque de $\varepsilon(Y)$ dans $C_X$
(\hyperref[I.4.4.1-fr]{I, 4.4.1}). La réciproque résulte aussi de
(\hyperref[II.8.2.12-fr]{8.2.12}) lorsque $Y$ est affine noethérien et $S$ de
type fini.

